\section{问题分析}\label{sec:analysis}

\subsection{缺失值与样本覆盖}
不同检验项目的缺失程度存在较大差异，且若干项目具有相同或接近的缺失模式。直接删除缺失比例较高的记录，会同时改变训练样本量和研究人群；若在评价时删除难以预测的记录，还会使所得误差偏离完整预测任务。因而，当前主协议保留全部训练、开发和B集记录，采用折内中位数填补与缺失指示，或采用模型原生缺失处理。

缺失模式可能包含检测安排或样本批次的信息，但现有材料不足以确定其形成机制。本文将缺失计数作为可验证的预测特征，不预设缺失为完全随机，也不赋予其未经证实的医学含义。

\subsection{相关性、量纲与特征表示}
多个体检指标存在统计相关，线性模型可能出现系数不稳定。正则化能够在保留相关变量的同时限制系数规模；标准化则改变不同量纲输入在模型训练中的相对尺度。PCA通过输入方差构造低维表示，其方差保留比例并不直接衡量血糖预测信息。因此，是否缩放、删列或压缩输入，应以训练内部预测比较为依据，而非仅依据相关图或累计方差曲线。

\subsection{非线性关系与误差互补}
线性模型为关联分析和预测提供基准，树模型能够表达轴向阈值与非线性交互，前馈网络则提供另一类连续函数表示。由于不同模型的归纳偏好和特征处理不同，其误差未必完全重合。融合是否有效，取决于同一批记录上的联合误差，而不能只由单模型排名决定。本文使用训练折外预测比较有限个融合配置，避免以模型对自身训练记录的拟合预测选择权重。

\subsection{评价边界与尾部误差}
MSE对较大误差更敏感，少量高目标值记录可能贡献较多平方损失。仅比较总体指标不足以说明误差分布，需要结合分箱诊断与配对误差差异分析。另一方面，表示选择、参数搜索和权重筛选均会引入选择偏差，重复使用开发或B集也会限制结论的外推。因此，本文分别报告训练内部选择、配置冻结后的已知集合复核及其条件不确定性。
